\documentclass[11pt]{article}
\usepackage{amsmath,amssymb,amsthm}
\newtheorem{theorem}{Theorem}
\newtheorem{proposition}{Proposition}
\begin{document}

\begin{theorem}[Arrangement-conditional pectinate quartet law]
Consider the hierarchical finite-state Moran model with species tree
$(((1,2),3),4)$, finite count processes on all branches, conditional binomial
inheritance at speciation, and the structured genealogy operators $A_2,A_3,A_4$
defined by the model. Let $Z=(Z_1,Z_2,Z_3,Z_4)$ be the sampled arrangement
pattern and let $Q_1=12|34$, $Q_2=13|24$, and $Q_3=14|23$. For each pattern
$z\in\{0,1\}^4$, construct pattern-specific terminal vectors
\[
u_{j,z_j}(k,b)=1_{\{b=z_j\}}p^{z_j}(1-p)^{1-z_j},
\qquad p=k/M.
\]
Propagate these vectors through the sister and intermediate recursions,
retaining the pattern mass, and let $F_z$ be the resulting unresolved
four-lineage root-entry vector. If $c_i(z)$ is the corresponding early
topology mass, then
\[
J[z,i]=P(Z=z,Q_i)=c_i(z)+
(\pi_{\rm old}\otimes h_i)^T e^{A_4T_A}F_z.
\]
Consequently,
\[
P(Z=z)=\sum_iJ[z,i],\qquad
P(Q_i)=\sum_zJ[z,i],\qquad
P(Q_i\mid Z=z)=\frac{J[z,i]}{\sum_jJ[z,j]}
\]
whenever $P(Z=z)>0$.
\end{theorem}

\begin{proof}
At a terminal count $k$, the event that the sampled arrangement is $z_j$ and
the sampled lineage carries background $b$ has probability
$1_{\{b=z_j\}}p^{z_j}(1-p)^{1-z_j}$. The two pattern vectors partition the
ordinary terminal sampling vector. Linearity of the dual one-lineage operator
and stochasticity of each inheritance kernel therefore preserve this
partition through every terminal and speciation stage.

At the intermediate node, unresolved pair states multiply conditionally on the
shared count. If the pair has already coalesced, the topology is $Q_1$; the
remaining taxon-3 pattern mass must still multiply that absorbing contribution.
The three-lineage operator partitions the resulting mass into its transient
states and the three absorbing topology states. The same argument applies at
the root: for early topology mass, the taxon-4 sampling probability is retained
conditional on the root count; for unresolved mass, $F_z$ contains the product
of the three-lineage and taxon-4 background values without normalization.

The residual vectors satisfy $h_1+h_2+h_3=1$ on every root state. Thus root
propagation partitions unresolved mass among $Q_1,Q_2,Q_3$. Summing over
topologies and patterns gives one, and summing in either order gives the two
marginals stated above. Division by the positive pattern marginal gives the
conditional law. If the pattern marginal is zero, no conditional probability
is defined and the implementation returns NaN.
\end{proof}

\begin{proposition}[A sufficient symmetry condition]
If the full conditional genealogy-and-sampling construction is invariant under
a taxon permutation that exchanges the two discordant quartet resolutions,
then $P(Q_2\mid Z=z)=P(Q_3\mid Z=z)$. In particular, a common fixed
background with no arrangement distinction reduces the two discordant
resolutions to exchangeable ordinary-coalescent outcomes.
\end{proposition}

\begin{proof}
The permutation gives a measure-preserving bijection between histories ending
in $Q_2$ and histories ending in $Q_3$, while preserving the specified pattern
event. Their joint masses are equal, and division by the common pattern mass
preserves equality.
\end{proof}

The theorem is exact for the hierarchical simulator model, including its
finite boundary rates and binomial daughter initialization. It is not a claim
about all biological Moran processes or about introgression.
\end{document}
